% Generated draft: requires mathematical and source review.
\documentclass[11pt]{article}
\usepackage[margin=25mm]{geometry}
\usepackage{fontspec}
\setmainfont{Latin Modern Roman}
\usepackage{amsmath,amssymb,mathtools,unicode-math}
\setmathfont{Latin Modern Math}
\usepackage{enumitem,needspace,xurl,hyperref}
\setlength{\parindent}{0pt}
\setlength{\parskip}{5pt}
\setlength{\emergencystretch}{3em}
\newcommand{\N}{\mathbb{N}}
\newcommand{\Z}{\mathbb{Z}}
\newcommand{\Q}{\mathbb{Q}}
\newcommand{\R}{\mathbb{R}}
\newcommand{\C}{\mathbb{C}}
\newcommand{\F}{\mathbb{F}}
\newcommand{\A}{\mathbb{A}}
\newcommand{\PP}{\mathbb{P}}
\newcommand{\E}{\mathbb{E}}
\newcommand{\eps}{\varepsilon}
\newcommand{\dd}{\,\mathrm{d}}
\begin{document}
\section*{UnsolvedMath ID: 2874}
\textbf{Source problem code:} catalogue entry
\subsection{Definitions and mathematical statement}
The source asks for the algebraic structure of several equivariant homology-cobordism groups, not one uniquely defined group. To explain the data without selecting an unsupported variant, an integral homology $3$-sphere is a closed oriented $3$-manifold $Y$ with $H_*(Y;\Z)\cong H_*(S^3;\Z)$. An ordinary smooth homology cobordism from $Y_0$ to $Y_1$ is a compact oriented smooth $4$-manifold $W$ with $\partial W=(-Y_0)\sqcup Y_1$ such that both inclusions induce integral homology isomorphisms. A rational version replaces $\Z$ by $\Q$.

One version described in the background uses pairs $(Y,\rho)$, where $Y$ is an integral homology sphere and $\rho:\Z/p\Z\to\operatorname{Diff}^+(Y)$ is an orientation-preserving action with nonempty fixed-point set. A schematic equivariant cobordism is a homology cobordism carrying a compatible action restricting to the boundary actions. This does not yet specify all equivalence conditions: smooth versus topological category, effectiveness, fixed-set behavior in the cobordism, allowed isotopies, and any basepoint or tangent-representation data for equivariant connected sum must be fixed. The source calls these ``appropriate'' classes but does not provide those full definitions, so no complete group presentation or operation is silently invented.

The notes distinguish more general versions recording a specified diffeomorphism on each homology sphere, without restricting the fixed-point set; a specified diffeomorphism need not be identified with a finite cyclic action. Integral and rational groups likewise remain distinct. For a chosen well-defined version, forgetting the extra action or diffeomorphism can give a homomorphism
\[
F:\Theta_{\mathrm{eq}}\longrightarrow\Theta,
\]
where $\Theta$ is the corresponding ordinary homology-cobordism group. These symbols are descriptive placeholders for a specified category, not a claim that all versions are canonically the same.

The requested structural programme includes torsion, free-abelian summands, other quotients and relations, kernels and cokernels of forgetful maps, and comparison of the fixed-point and diffeomorphism variants. There is no single conjectural group decomposition in the literal question.

The supplied literature triage dated 2026-08-17 reports Dai--Hedden--Mallick's involutive refinement with an infinite-rank free-abelian subgroup generated by strongly nonextendable cork boundaries, including the $p=2$ setting. Here $\Z^\infty$ means the direct sum $\bigoplus_{j\geq1}\Z$, not an unrestricted direct product. The source reports that a 2025 Hendricks--Mallick--Stoffregen--Zemke theorem gives a $\Z^\infty$ direct summand in a forgetful-map kernel for a particular equivariant group. That stronger summand statement is not weakened to mere subgroup existence, nor extended to every cyclic-action variant or every prime.

The background also records suitable branched-cover homomorphisms from ordinary smooth knot concordance and strongly invertible concordance into equivariant rational homology-cobordism groups, with the branching action retained, and related equivariant Floer obstructions to sliceness and concordance. The exact cover orders, group definitions, and map hypotheses are not supplied, so no map into an integral group or arbitrary-cover theorem is inferred.

The registry is open and the dated triage classifies progress as partial: large subgroups and kernels are known, but there is no overall classification of torsion, quotients, relations, or category dependence. All theorem, map, and status reports are supplied information rather than independently verified results in this draft.
\subsection{Short English statement}
Determine the structure of equivariant homology-cobordism groups, including torsion, summands, forgetful-map kernels and cokernels, and comparisons between variants. The supplied literature detects infinite-rank subgroups and a kernel summand, but the action, fixed-set, coefficient, and equivalence conventions must first specify the group being studied.
\subsection{Sources}
\begin{itemize}
\item R. Inanc Baykur, Robion C. Kirby, and Daniel Ruberman, editors, K3: A New Problem List in Low-Dimensional Topology, AMS Mathematical Surveys and Monographs 295 (2026), Problem 3.76: \url{https://math.berkeley.edu/sites/default/files/surv-295-ruberman-watermarked-author-pdf.pdf}. This citation is retained without opening or generating a PDF.
\item [DHM23] Irving Dai, Matthew Hedden, and Abhishek Mallick, ``Corks, involutions, and Heegaard Floer homology,'' Journal of the European Mathematical Society 25(6) (2023), 2319--2389: \url{https://doi.org/10.4171/jems/1239}.
\item [AB24b] Antonio Alfieri and Keegan Boyle, ``Strongly invertible knots, invariant surfaces, and the Atiyah-Singer signature theorem,'' Michigan Mathematical Journal 74(4) (2024), 845--861, DOI 10.1307/mmj/20226183.
\item [BH24a] David Baraglia and Pedram Hekmati, ``Equivariant Seiberg-Witten-Floer cohomology,'' Algebraic and Geometric Topology 24(1) (2024), 493--554, DOI 10.2140/agt.2024.24.493.
\item Kristen Hendricks, Abhishek Mallick, Matthew Stoffregen, and Ian Zemke, ``The link surgery formula and equivariant surgeries,'' arXiv:2507.12809 (2025): \url{https://arxiv.org/abs/2507.12809}.
\item Local registry data/unsolvedmath/problems.json, permanent catalogue entry, source code catalogue entry, ``Kirby Problem 3.76''; statement, complete supplied literature notes, and literature triage dated 2026-08-17. Dataset attribution: Ulam AI and the UnsolvedMath Contributors.
\end{itemize}
\end{document}
